% v6 draft section (Claude, cc/iina-tier1). Every number is a generated \vsix macro (scripts/hero_v6.py,
% scripts/v6_macros.py); none is typed here. To be \input during the post-PR-6 integration pass.
\section{Stopping a Model Is Not Stopping Its Effects}\label{sec:v6stop}

An operator who decides to stop an agent can act at three places. \emph{Halting inference} ends the stream of new
proposals. \emph{Revoking admission} at the broker refuses permits for proposals that arrive later. \emph{Fencing the
sink} raises the epoch floor at the component that writes durable state, so that any operation carrying an older epoch
is refused at commit, including operations that were admitted, dispatched and are already running. The first two are
what most deployed stop controls do~\cite{controlroom}. Only the third reaches work that is already in flight, and it
is the fencing-token argument of distributed locking~\cite{kleppmann2016} applied to agent authority.

\para{Setup (E4).} We run the real broker and sink processes (Rust, SQLite, JSON-lines RPC; Section~\ref{sec:model})
under a synthetic agent that emits \(\lambda = 2\) tool calls per second to \(C \in \{4, 32\}\) concurrent executors
whose latencies are exponential with mean \(w \in \{0.5, 5, 30\}\)\,s. Each call is a pre-approved, authorized one-unit
effect; nothing in the workload is adversarial. After a 20\,s warm-up the operator decides to stop, and the stop takes
effect \(L \in \{0, 1, 5, 35\}\)\,s later. A trial counts the effects the sink commits after the stop took effect,
read from the sink's own database, and the legitimate calls the stop discards (cancelled or denied). The plan, grid and
metrics were committed before the first trial (\texttt{fixtures/v6/revokebench\_plan.json}); the \vsixEFourTrials{}
trials include \vsixEFourFaultTrials{} with SIGKILL of the sink or delayed completions.

\para{Result.} Figure~\ref{fig:v6hero}a shows the stable cell with the longest latency (\(C = \vsixStopCalls\),
\(w = \vsixStopLatency\)\,s, \(L = 0\)), chosen by that rule rather than by eye. After the stop took effect, halting
inference let \vsixHaltRange{} effects commit and revoking at the broker \vsixRevokeRange{}, the last
\vsixLastCommitSeconds\,s later; fencing the sink let \vsixFenceRange{} commit and discarded \vsixFenceLost{}
legitimate calls. Across all \vsixEFourCells{} fault-free cells, halting leaked up to \vsixGridHaltMax{} effects and
revoking up to \vsixGridRevokeMax; the fence admitted at most \vsixFencePostMax{} in \vsixFenceTrials{} trials.

\para{What the numbers mean.} The fence's zero is a property of the design, not a discovery: the sink serializes its
RPCs, \texttt{fence} cancels every pending record below the new floor, and effect identifiers are unique in the schema,
so duplicates are impossible rather than rare. What E4 measures is (i) that this implementation keeps the property
under the faults we injected, and (ii) the size and cost of the window the other two layers leave open. That window is
queueing, and it behaves like queueing. Where the offered load \(\rho = \lambda w / C\) is below one, halting leaks
roughly the calls in flight, \(\approx \lambda w\) (Little's law), and revocation leaks the calls already accepted.
Where \(\rho \ge 1\) (\vsixEFourOverloadedCells{} of the \vsixEFourWCCells{} \((w, C)\) cells), the backlog grows during the warm-up and
halting leaks all of it; we report those cells separately and do not apply Little's law to them. The fence pays for its
zero in the opposite currency. Every legitimate call in flight at the stop is discarded, which is what an operator who
stops an agent usually wants and what an operator who meant ``pause'' does not.

\para{Faults (E4b).} A SIGKILL that lands between requests tests little. In a separately predeclared run
(\texttt{fixtures/v6/revokebench\_b\_plan.json}) we kill the sink while a \texttt{complete} or a \texttt{fence} request is
outstanding, hold the executor's \texttt{accept} across the fence, and drive the system with closed-loop agents that
wait for each result. \vsixEFourBSentence

\para{Scope.} The sink mediates only its own SQLite effects. An external effect (a payment API, a repository, a cloud
resource) is fenced only if that surface can refuse a stale epoch, which is the question of
Section~\ref{sec:v6boundary}. The latencies are sampled, not production traces, and the agent is synthetic.
